
\begin{mechanism}[title=Beschreiben Sie die Umsetzung der genetischen Information in Eukaryoten unter Benennung der relevanten Zwischenstufen.]
\begin{enumerate}
    \item Nukleus: Transkription
    \item Nukleus: Bildung der pre-mRNA durch Anfügen eines CAPs am 5' Ende.
    \item Nukleus: Bildung der polyadenylated
pre-mRNA durch Anfügen einer poly A Kette am 3' Ende.
    \item Nukleus: Bildung der mRNA durch splicing.
    \item Cytoplasma: Translation der mRNA
    \item ER und Golgi-Apparat: Faltung der RNA in die unterschiedlichen Strukturen bis zum finalen Protein.
\end{enumerate}
\end{mechanism}
\begin{mechanism}[title=Definition Nukleotid]
Molekül bestehend aus: 
\begin{itemize}
    \item \textbf{Zucker} (Desoxyribose/Ribose)
    \item Eine oder mehrere \textbf{Phosphatgruppen}
    \item Eine Stickstoffhaltige \textbf{Base}
\end{itemize}
\end{mechanism}
\begin{mechanism}[title=]

\end{mechanism}
\begin{mechanism}[title=]

\end{mechanism}
